% q0025.tex — Maximizing shareholder value [Numeric]
\question{
  id        = {0025},
  title     = {Maximizing Shareholder Value},
  source    = {OBECON},
  year      = {2026},
  round     = {Seletiva IEO -- Phase A},
  language  = {English},
  topic     = {Finance},
  subtopic  = {Corporate Finance / IRR},
  type      = {numeric},
  statement = {%
    A manufacturing company evaluates an M\&A acquisition costing \$100M
    upfront. Short-term CF scenarios: Base (45\%, \$80M), Bear (40\%,
    half of base), Bull (15\%, four times base). Long-term CF: \$200M
    certain. Shareholders are risk-neutral.

    Calculate the IRR per period. Give as \% rounded to nearest integer;
    if 0.5, round to even.
  },
  choices   = {},
  answer    = {100},
  solution  = {%
    \textbf{Step 1 --- Expected short-term CF:}
    \[
      E[\text{CF}_1] = 0.45 \times 80 + 0.40 \times 40 + 0.15 \times 320
      = 36 + 16 + 48 = \$100\text{M}
    \]

    \textbf{Step 2 --- Set NPV = 0:}
    \[
      -100 + \frac{100}{1+r} + \frac{200}{(1+r)^2} = 0
    \]
    Let $x = \dfrac{1}{1+r}$:
    \[
      200x^2 + 100x - 100 = 0 \implies 2x^2 + x - 1 = 0
      \implies (2x - 1)(x + 1) = 0
    \]
    Taking the positive root: $x = \dfrac{1}{2}$, so $1 + r = 2$,
    i.e.\ $r = 1 = \mathbf{100\%}$ per period.
  },
}
